\section{Analisi dei dati}

In questo capitolo vengono presentati i risultati ottenuti mediante l'analisi dei dati (quantitativi e qualitativi) raccolti durante gli esperimenti con gli utenti.

Le due versioni del sito (la versione originale e la versione post-riprogettazione) sono state testate con l'obiettivo di confrontarle e verificare l'effettiva efficacia delle modifiche implementate. Inoltre, l'esperimento è servito per individuare problemi nella versione originale che non erano stati rilevati durante la valutazione euristica ed eventuali criticità residue nella versione post-riprogettazione.

Per l'esecuzione dell'esperimento è stato scelto l'approccio \textit{within-subjects} (a misure ripetute), in cui ogni utente ha testato entrambe le versioni del sistema. Questa modalità è stata ritenuta la più adeguata per massimizzare la validità statistica delle misurazioni e annullare la variabilità dovuta alle differenze individuali dei partecipanti. L'esperimento ha coinvolto un totale di 15 utenti.

Per limitare l'effetto di apprendimento derivante dall'utilizzo di entrambe le versioni da parte dello stesso utente, è stato applicato un bilanciamento dell'ordine di esecuzione. Nello specifico, il campione è stato suddiviso in questo modo:
\begin{itemize}
	\item 9 utenti hanno testato prima la versione originale e successivamente quella post-riprogettazione (ordine A-B);
	\item 6 utenti hanno testato prima la versione post-riprogettazione e successivamente quella originale (ordine B-A).
\end{itemize}

Tuttavia, si è scelto di non effettuare una suddivisione perfettamente equa, assegnando un peso maggiore alla sequenza A-B. Tale decisione è stata motivata dalla volontà di simulare più fedelmente uno scenario reale di aggiornamento del sistema: l'obiettivo è stato quello di raccogliere un maggior numero di osservazioni sul comportamento tipico di un utente che passa dal sistema originale a quello riprogettato. Il sottogruppo B-A è stato comunque mantenuto per garantire il rigore statistico necessario a isolare e valutare l'effetto di apprendimento nell'analisi dei dati.

Nelle sezioni successive verranno dapprima presentate le caratteristiche del campione di partecipanti, per poi passare all'analisi dei dati raccolti per i singoli compiti. Infine, verranno analizzati i risultati del questionario SUS (\textit{System Usability Scale}) e i commenti raccolti durante l'intervista finale.
